\section{Strict two-categorical indexing: the ordinary-base case}

\paragraph{J.1 (Proved).}
Let $S$ be an ordinary small category, regarded as a locally
discrete strict $2$-category. On $2\operatorname{Cat}/S$ define
the presheaf fine class by requiring $s\backslash u$ to be a
Street weak equivalence for every $s\in S$. Its infinity-localization
is equivalent to $\operatorname{Fun}(NS^{\mathrm{op}},\mathcal S)$.
The dual over-slice class gives
$\operatorname{Fun}(NS,\mathcal S)$. These equivalences are
compatible with the Street-nerve comparisons to the contravariant
and covariant simplicial theories over $NS$.
\begin{proof}
Because the base has only identity $2$-cells, its lax under-slice
agrees with the strict $2$-comma object. Define
\[
 G(X\xrightarrow p S)(s)=s\backslash X,
 \qquad G(X):S^{\mathrm{op}}\longrightarrow2\operatorname{Cat},
\]
with strictly functorial reindexing by precomposition. For a
strict diagram $A:S^{\mathrm{op}}\to2\operatorname{Cat}$, define
its $2$-categorical Grothendieck construction $E(A)$ as follows.
Objects are $(s,x)$ with $x\in A(s)$; a $1$-cell
$(s,x)\to(t,y)$ is $(a:s\to t,h:x\to A(a)y)$; a $2$-cell
between such arrows exists only when their base arrows agree,
and is the corresponding $2$-cell between their $h$ components.
Strict composition in $A$ gives strict composition in $E(A)$.

For every $s$ there is a strict $2$-functor
\[
 R_s:s\backslash E(A)\longrightarrow A(s),
 \quad((t,y),a:s\to t)\longmapsto A(a)y.
\]
Its section takes $x$ to $((s,x),1_s)$. The arrows $(a,1)$
give a strict $2$-natural transformation from section followed
by $R_s$ to the identity. Thus $R_s$ is a Street weak
equivalence; a strict $2$-natural transformation supplies a
simplicial homotopy after the usual cylinder comparison. The
$R_s$ assemble naturally to $GE(A)\to A$.

There is also a natural strict $2$-functor over $S$,
\[
 j_X:X\longrightarrow EG(X),\qquad
 x\longmapsto(p(x),(x,1_{p(x)})).
\]
At $s$, the composite $R_s(s\backslash j_X)$ is the identity
on $s\backslash X$. For an object
$((t,(x,\alpha:t\to p(x))),\beta:s\to t)$ of
$s\backslash EG(X)$, the arrow with base $\alpha$ and
identity component gives a natural transformation to its image
under $(s\backslash j_X)R_s$. Consequently
$s\backslash j_X$ is a Street weak equivalence. Notice that
all displayed equalities are strict because the base is an
ordinary category; this is the precise place where that
hypothesis enters.

By definition $G$ preserves and detects the fine class. The
natural equivalence $GE\to1$ proves that $E$ sends pointwise
Street weak equivalences to fine equivalences. Together with
$j:1\to EG$, these are inverse equivalences after infinity-localization.
The projective model structure on strict $S^{\mathrm{op}}$-diagrams
in the combinatorial Thomason model category of strict
$2$-categories presents
$\operatorname{Fun}(NS^{\mathrm{op}},\mathcal S)$: apply the
projective diagram Quillen equivalence induced by
$c_2\operatorname{Sd}^2\dashv\operatorname{Ex}^2N_2$
\citep[Theorem 6.27 and Corollary 6.32]{GHIJAraMaltsiniotis2014}.
The standard strict-diagram comparison is
\citet[Proposition 4.2.4.4]{GHIJLurie2009}, applied on the
simplicial-set side of that Quillen equivalence.

To identify the comparison functor, observe the strict identity
\[
 N_2(s\backslash X)
    =N(S_{s/})\times_{NS}N_2X.
\]
It follows because the Street nerve preserves limits and the
base comma construction is the indicated strict pullback.
In the contravariant straightening theory over an ordinary
nerve, the value at $s$ is the weak homotopy type of this
pullback with the representable left fibration $N(S_{s/})$.
This is the co-Yoneda calculation for straightening: on a
simplex $\Delta^n\to NS$ the pullback has one contractible
component for each arrow $s\to\sigma(n)$, and the general
calculation follows by the simplicial resolution and homotopy
colimits. Thus the Street nerve has the same straightened
diagram as the equivalence already constructed. The over-slice
case follows by total $2$-categorical duality and invariance of
Street weak equivalences under that duality \citep[Theorem 6.6]{GHIJAraMaltsiniotis2022}.
\end{proof}

\paragraph{J.2 (Open).}
The proof of J.1 does not establish the requested recognition
for the walking noninvertible or invertible $2$-cell, the
one-object generating-$2$-cell examples, or a general strict
$2$-category. For those bases the strict comma identity used
above fails: lax triangles contain nonidentity $2$-cells, and
reindexing must retain their coherent action. No relative
strictification or straightening theorem with all the required
hypotheses has been proved here.

\paragraph{J.3 (Open).}
The universally pullback-preserving classes in Part J(2), the
complete boundary-piece tests, and the proposed lax-comma Der~4
base-change theorem in J(3) are not established here. The
absolute model structure and Theorem A do not imply those
claims. The verified statements available from the sources
are Theorem A for the specified slices
\citep[Theorem 5.2]{GHIJAraMaltsiniotis2018} and invariance
under every cell duality
\citep[Theorem 6.6]{GHIJAraMaltsiniotis2022}. Reversing only $2$-cells exchanges lax
and oplax coherence; reversing all positive-dimensional cells
exchanges the matched under/over fine theories. A full
four-variance base-change table requires the still-missing
recognition and lax-comma compatibility.
